\chapter{Background and Literature Review}\label{ch:background}

This chapter reviews the literature that the thesis builds on and identifies the gaps it addresses. It is organised by idea rather than by author, following the perceive--reason--act loop of Chapter~\ref{ch:intro}. Section~\ref{sec:bg:shift} states the common problem: distribution shift, compounding error and the long tail. Sections~\ref{sec:bg:repr} to~\ref{sec:bg:manip} review the literature on perception (3D scene representations), on reasoning (language and foundation models in driving), on acting in driving (end-to-end planning and its evaluation), on adaptation to shift, and on robot manipulation. Section~\ref{sec:bg:bench} reviews benchmarks and challenges as research infrastructure, together with the statistics needed to interpret their results, which is the subject of the methodology of Chapter~\ref{ch:methodology}. Each section closes by stating what is known and what is not, and Section~\ref{sec:bg:summary} collects the resulting gaps and relates them to the research questions. Works led by the author are discussed in the chapters based on them. Here they, and works that the author co-authored, are cited only to locate the thesis within each line of work, and co-authored works are marked as such.

% ==============================================================================
\section{Distribution Shift for Physical Agents}\label{sec:bg:shift}

\subsection{Covariate shift and compounding error}

In the classical formulation, covariate shift changes the distribution of inputs while the conditional distribution of outputs stays fixed, and can be corrected by importance-weighting the training likelihood~\cite{shimodaira2000improving}. Sequential decision-making breaks the assumption that makes this correction possible: the learner's own actions generate the states on which it is later tested. Behaviour cloning therefore suffers from compounding errors, with a regret that grows quadratically in the task horizon~\cite{ross2010efficient}. DAgger reduces imitation learning to no-regret online learning by aggregating expert labels on the states that the learner itself visits, so that the resulting policy performs well under the distribution it induces~\cite{ross2011dagger}.

Because on-policy querying requires an interactive expert, later work widened the training distribution offline instead. DART injects noise into the demonstrator during data collection~\cite{laskey2017dart}, and ChauffeurNet synthesises perturbations of expert driving after finding that behaviour cloning was insufficient for complex driving even with 30 million examples~\cite{bansal2019chauffeurnet}. Human-gated schemes let the expert take control only when needed~\cite{kelly2019hgdagger}. The debate is therefore about where the corrective signal comes from: an interactive expert, which is accurate but costly and sometimes unsafe; designer-specified perturbations, which are cheap but limited to anticipated deviations; or the deployed system itself, which is the option explored in Chapters~\ref{ch:centaur} and~\ref{ch:kai0}.

\subsection{Causal confusion and shortcut learning}

A behaviour-cloning policy is a discriminative model that ignores the causal structure linking the expert, the environment and the action. Under distribution shift this produces causal misidentification, in which access to more information can yield worse performance~\cite{dehaan2019causal}. A frequent special case is the copycat problem, where an imitator with access to observation history learns to repeat the previous action rather than to predict the next one~\cite{wen2020fighting}. In driving, dataset bias, causal confusion and training variance have been identified as limits of behaviour cloning~\cite{codevilla2019exploring}, and the wider machine-learning literature describes the same phenomenon as shortcut learning: decision rules that score well on benchmarks but fail to transfer~\cite{geirhos2020shortcut}. Open-loop evaluation of driving shows the problem concretely: on nuScenes, planners that use only the ego vehicle's own state are competitive with full perception-based stacks~\cite{zhai2023rethinking,li2024ego}. Proposed remedies include targeted interventions or expert queries~\cite{dehaan2019causal}, removal of nuisance information from the representation~\cite{wen2020fighting}, and privileged or structured intermediate representations that make the relevant causes explicit~\cite{chen2019lbc}. Whether large-scale pre-training mitigates or amplifies causal confusion remains open, and recent reinforcement-learning approaches to driving still cite it as their motivation~\cite{gao2025rad}.

\subsection{The long tail and the statistics of rare events}

Demonstrating the reliability of an autonomous vehicle by accumulated test mileage is statistically prohibitive~\cite{kalra2016driving}. Liu and Feng call this the curse of rarity: safety-critical events are rare and embedded in high-dimensional driving scenes, so both learning and validation become data-inefficient~\cite{liu2024curse}. Their dense reinforcement-learning approach accelerates validation by concentrating on safety-critical states~\cite{feng2023dense}. The long tail is not confined to planning: corner-case datasets expose detection failures on unusual road objects~\cite{li2022coda}, and cross-dataset studies of 3D detection trace much of the adaptation gap to differences in vehicle sizes between regions~\cite{wang2020train}. Three responses compete: collecting more data, generating and testing rare situations in a targeted way, and equipping the deployed system with mechanisms to cope with them. Robustness to distributional shift was listed among five concrete problems of AI safety a decade ago~\cite{amodei2016concrete}, yet research on physical agents still lacks an agreed notion of how much of the relevant long tail a dataset or benchmark covers.

\subsection{Detecting shift: out-of-distribution detection and uncertainty}

A simple baseline scores test inputs by their maximum softmax probability~\cite{hendrycks2017baseline}. Approximate Bayesian methods estimate predictive uncertainty through Monte-Carlo dropout~\cite{gal2016dropout} or deep ensembles, which also assign higher uncertainty to inputs outside the training distribution~\cite{lakshminarayanan2017simple}, and aleatoric uncertainty can be separated from epistemic uncertainty~\cite{kendall2017uncertainties}. A generalised framework unifies anomaly, novelty, open-set and outlier detection~\cite{yang2024generalized}, and corruption benchmarks measure robustness to common perturbations~\cite{hendrycks2019benchmarking}. In driving, robust imitative planning uses epistemic uncertainty to detect unfamiliar scenes and to adapt with a few expert demonstrations~\cite{filos2020can}. Two limitations matter for this thesis. The quality of uncertainty estimates itself degrades as shift increases~\cite{ovadia2019can}, so detectors calibrated on in-distribution data can fail exactly where they are needed. And most of this literature studies classification: for planners that legitimately output several good trajectories, naive predictive entropy conflates valid multimodality with ignorance, a problem that the Cluster Entropy of Chapter~\ref{ch:centaur} addresses.

\paragraph{What is known and what is not.} It is established that imitation errors compound with the horizon and that on-policy data correct this in principle~\cite{ross2010efficient,ross2011dagger}; that richer inputs can hurt through causal confusion~\cite{dehaan2019causal,codevilla2019exploring}; that rare events make validation by exposure infeasible~\cite{kalra2016driving,liu2024curse}; and that uncertainty estimates degrade under shift~\cite{ovadia2019can}. What is not known is how a deployed physical agent can obtain an on-policy corrective signal without an expert in the loop, and how uncertainty can be turned from a passive alarm into a learning signal for multimodal planners and manipulation policies. These gaps motivate RQ3 and RQ4.

% ==============================================================================
\section{3D Scene Representations for Driving}\label{sec:bg:repr}

\subsection{Object-centric 3D boxes}

The dominant interface between perception and planning has long been a list of 3D bounding boxes. LiDAR detectors encode point clouds into vertical pillars for speed~\cite{lang2019pointpillars} or detect object centres, which also simplifies tracking~\cite{yin2021centerpoint}. Camera-based detectors range from fully convolutional monocular detection~\cite{wang2021fcos3d} to multi-view, query-based detection that links 3D object queries to image features by projection~\cite{wang2021detr3d}. Boxes are compact and match the interfaces of downstream prediction and planning, but they presuppose a closed vocabulary of categories and a cuboid shape. Box-centric methods neglect finer geometric detail and struggle with general, out-of-vocabulary objects~\cite{tian2023occ3d}, and corner-case datasets confirm that unusual objects are a real source of failure~\cite{li2022coda}.

\subsection{Bird's-eye view as a shared plane}

Lift-Splat-Shoot lifts image features into 3D using predicted depth distributions and splats them onto a bird's-eye-view (BEV) grid~\cite{philion2020lss}. BEVDet built a multi-camera detector on this plane~\cite{huang2021bevdet}, BEVDepth added explicit depth supervision~\cite{li2023bevdepth}, and BEVFusion fused camera and LiDAR features in one BEV space for several tasks~\cite{liu2023bevfusion}. BEVFormer (co-authored by the author), which gathers BEV features from multiple cameras and time steps with spatiotemporal attention, belongs to this line~\cite{li2022bevformer}, as does a survey of BEV perception that the author also co-authored~\cite{li2023bevdevils}. BEV collapses the height axis, which loses overhanging and vertically complex structure, and camera-only BEV quality is bounded by depth estimation. Its lasting role has been as the intermediate feature space of many end-to-end planners (Section~\ref{sec:bg:e2e}).

\subsection{Vectorised maps}

HDMapNet reframed high-definition maps as something constructed online from onboard sensors~\cite{li2022hdmapnet}. MapTR models each map element as a point set with a group of equivalent permutations, which stabilises the learning of vectorised maps~\cite{liao2023maptr}, and vectorised scene representations can feed planning directly~\cite{jiang2023vad}. OpenLane-V2 (co-authored by the author) extends this line from lane lines to lane topology and the relations between lanes and traffic elements~\cite{wang2023openlanev2}. Vector maps are compact and planning-friendly, but only for predefined element types; unstructured areas and arbitrary obstacles fall outside their vocabulary.

\subsection{Semantic occupancy}

Occupancy grids date back to the probabilistic, tessellated maps of mobile-robot navigation~\cite{elfes1989occupancy}. Semantic scene completion jointly infers geometry and semantics from a single depth image~\cite{song2017sscnet}, and MonoScene brought it to monocular RGB~\cite{cao2022monoscene} on the outdoor SemanticKITTI benchmark~\cite{behley2019semantickitti}. Surround-view methods followed, including the tri-perspective view of TPVFormer~\cite{huang2023tpvformer} and SurroundOcc, which lifts multi-scale features with 2D--3D attention and builds dense labels by fusing multi-frame LiDAR~\cite{wei2023surroundocc}. Benchmarks appeared at the same time: Occ3D, with a visibility-aware label-generation pipeline on Waymo and nuScenes~\cite{tian2023occ3d}, and OpenOccupancy~\cite{wang2023openoccupancy}. OccNet and its OpenOcc benchmark, of which the author is co-first author~\cite{sima2023occnet}, and the OpenScene dataset release~\cite{sima2023openscene} belong to this line and are the subject of Chapter~\ref{ch:occnet}.\footnote{Despite the similar name, OccNet is unrelated to the Occupancy Networks of Mescheder et al.~\cite{mescheder2019occupancy}, which represent the surface of a single object implicitly, as the decision boundary of a classifier over continuous 3D coordinates.}

The main debate within this line concerns efficiency versus completeness. Dense volumes ignore the sparsity of scenes and are computationally expensive. FB-OCC, the winning entry of the CVPR 2023 occupancy challenge, combined forward and backward view transformation with joint depth and semantic pre-training, model scaling and post-processing~\cite{li2023fbocc}. Later designs sought comparable accuracy at lower cost. FlashOcc keeps features in the BEV plane and lifts only its output to 3D by a channel-to-height transformation~\cite{yu2023flashocc}, and sparse methods compute only where the scene is occupied: as sparse latent features~\cite{tang2024sparseocc}, as sparse voxels processed by a fully sparse network~\cite{liu2024sparseocc}, or as a set of points and classes predicted directly by a transformer~\cite{wang2024opus}. Object-like 3D Gaussians offer a middle ground~\cite{huang2024gaussianformer,huang2025gaussianformer2}, borrowing the primitive of Gaussian-splatting rendering~\cite{kerbl2023gaussian}. The common finding is that dense voxel features, as used by OccNet in Chapter~\ref{ch:occnet}, are not needed for competitive accuracy on these benchmarks.

A second, less discussed issue is that occupancy labels are generated automatically from aggregated LiDAR under different conventions (visibility masks, voxel sizes, class sets), so scores are not comparable across benchmarks and label noise is entangled with model error. The metric has changed as well. Voxel-level mIoU penalises errors inconsistently along the depth axis. RayIoU instead casts query rays into the predicted volume, as a LiDAR sensor would, and counts a ray as correct only if the first occupied voxel it meets has the right class and lies within a distance threshold of the true depth~\cite{liu2024sparseocc}. The two public occupancy challenges reflect both issues.\footnote{Challenge pages: \url{https://opendrivelab.com/challenge2023/} and \url{https://opendrivelab.com/challenge2024/}.} The 3D Occupancy Prediction Challenge at CVPR 2023 used the Occ3D-nuScenes labels (0.4\,m voxels, with camera-visibility masks) and ranked entries by voxel mIoU. The Occupancy and Flow track of the CVPR 2024 Autonomous Grand Challenge used OpenOcc v2, a revised version of the benchmark of Chapter~\ref{ch:occnet}, and scored entries by RayIoU together with the mean absolute velocity error (mAVE) of occupancy flow. The results of Chapter~\ref{ch:occnet}, reported as voxel mIoU on the first version of OpenOcc, are therefore not directly comparable with later results on either benchmark.

Two further lines address limitations that Chapter~\ref{ch:occnet} leaves open. The first reduces the dependence on 3D labels. RenderOcc renders depth and semantic maps from a NeRF-style volume and supervises them with 2D labels only~\cite{pan2024renderocc}. SelfOcc learns occupancy from video alone, by rendering neighbouring frames from a signed-distance representation~\cite{huang2024selfocc}, and GaussTR aligns splatted 3D Gaussians with the features of 2D foundation models~\cite{jiang2025gausstr}. The second line opens the vocabulary. POP-3D learns a voxel map of language-aligned embeddings from images, LiDAR and a pre-trained vision-language model, so that the scene can be queried with free-form text~\cite{vobecky2023pop3d}. VEON adapts a depth model and a vision-language model to lift open-vocabulary semantics into 3D~\cite{zheng2024veon}, and LangOcc distils CLIP features into a voxel grid through volume rendering, from camera images only~\cite{boeder2025langocc}. Both lines are promising, but the gap in accuracy remains large. Models trained without manual semantic labels reach about 12--15 mIoU on Occ3D-nuScenes~\cite{jiang2025gausstr,zheng2024veon}, whereas the supervised winner of the 2023 challenge reached 54.19 mIoU on its test split~\cite{li2023fbocc}, although the two settings also differ in model size and evaluation split. Open-vocabulary occupancy is, moreover, still evaluated mainly on the closed classes of existing benchmarks, supplemented by small sets of free-form queries~\cite{vobecky2023pop3d}, so how much of the long tail it covers has not been measured.

Occupancy has also become a target for self-supervised pre-training. UniPAD pre-trains camera and LiDAR backbones by 3D volumetric differentiable rendering and reports gains in 3D detection and segmentation on nuScenes~\cite{yang2024unipad}. SPOT shows that occupancy prediction learns LiDAR representations that transfer across datasets, sensors and tasks, and that downstream performance improves with more pre-training data~\cite{yan2025spot}. Both evaluate transfer to perception tasks.

\subsection{Occupancy as the state of a world model}

Occupancy has also become the state space of world models. Cam4DOcc established a benchmark for camera-only 4D occupancy forecasting on nuScenes and Lyft data~\cite{ma2024cam4docc}. OccWorld forecasts future scenes and ego motion in 3D occupancy space~\cite{zheng2024occworld}, OccSora generates 4D occupancy as a driving world simulator~\cite{wang2026occsora}, OccLLaMA places occupancy, language and action tokens in one autoregressive model~\cite{wei2024occllama}, and UniScene first generates semantic occupancy and then video and LiDAR conditioned on it~\cite{li2025uniscene}. Copilot4D learns unsupervised world models of LiDAR observations with discrete diffusion~\cite{zhang2024copilot4d}, and latent world models have also served as self-supervision for end-to-end planning~\cite{li2025law}. Explicit occupancy world models are interpretable and checkable, whereas latent ones are cheaper and avoid dependence on labels. Most occupancy world models are judged by the quality of the occupancy they forecast or generate; where they also plan, as OccWorld and OccLLaMA do, planning is evaluated open-loop on nuScenes~\cite{zheng2024occworld,wei2024occllama}, the protocol whose weaknesses are discussed in Section~\ref{sec:bg:e2e}. How to evaluate a world model by its usefulness for planning rather than by reconstruction metrics therefore remains open.

\paragraph{What is known and what is not.} Dense 3D occupancy captures geometry beyond fixed vocabularies, can be predicted from cameras, can be made sparse or Gaussian for efficiency, and can serve as the state of forecasting models~\cite{wei2023surroundocc,tian2023occ3d,liu2024sparseocc,huang2024gaussianformer,zheng2024occworld}. It can also be learned without 3D labels or queried with an open vocabulary, so far at a large cost in accuracy~\cite{huang2024selfocc,vobecky2023pop3d,jiang2025gausstr}, and its evaluation is moving from voxel mIoU to ray-based metrics~\cite{liu2024sparseocc}. What is not known is whether occupancy makes downstream planning better than boxes or vectors beyond open-loop comparisons, and whether a single representation can serve detection, segmentation, forecasting and planning at once, at a scale that matches the diversity of real driving. This is the gap behind RQ1, which Chapter~\ref{ch:occnet} addresses, for planning with open-loop evidence only.

% ==============================================================================
\section{Language and Foundation Models in Driving}\label{sec:bg:lang}

\subsection{Datasets for explanation, grounding and question answering}

Language entered driving through explanation, with the textual justifications of BDD-X~\cite{kim2018bddx} and the explanations attached to action-inducing objects in BDD-OIA~\cite{xu2020bddoia}, and through grounding, with the object-referring commands of Talk2Car~\cite{deruyttere2019talk2car}. Question answering followed: NuScenes-QA poses questions over images and point clouds~\cite{qian2024nuscenesqa}, NuPrompt attaches language prompts to 3D perception~\cite{wu2025nuprompt}, Reason2Drive targets chain-based reasoning~\cite{nie2024reason2drive}, and LingoQA provides 28K short video scenarios with 419K annotations, on which GPT-4V answered truthfully on 59.6\% of questions against 96.6\% for humans~\cite{marcu2024lingoqa}. OmniDrive generates counterfactual reasoning data to align agents with 3D driving tasks~\cite{wang2025omnidrive}, and DriveLMM-o1 annotates perception, prediction and planning questions with step-by-step reasoning and scores the reasoning steps as well as the final answer~\cite{ishaq2025drivelmmo1}. In general vision, compositional questions generated from scene graphs were introduced by GQA~\cite{hudson2019gqa}. Most driving question-answering datasets, however, treat question--answer pairs as independent, so they cannot test whether an answer about planning is logically supported by the answers about perception and prediction on which it depends. DriveLM, of which the author is co-first author and which is the subject of Chapter~\ref{ch:drivelm}, addresses this point~\cite{sima2024drivelm}.

\subsection{Question decomposition and compositional reasoning}

Graph-structured question answering draws on work that decomposes a complex question into simpler ones. In language, the question decomposition meaning representation writes a question as the ordered list of steps needed to answer it, collected at scale in the Break benchmark~\cite{wolfson2020break}, and least-to-most prompting has a language model solve a sequence of sub-problems, each with the help of the answers to earlier ones~\cite{zhou2023leasttomost}. Tree- and graph-structured prompting generalise such chains to searches over intermediate thoughts~\cite{yao2023tree,besta2023got}. In vision, VisProg and ViperGPT let a language model compose calls to vision modules into a program whose intermediate outputs form an inspectable rationale~\cite{gupta2023visprog,suris2023vipergpt}, and Visual CoT grounds intermediate reasoning steps in image regions~\cite{shao2024visualcot}. Decomposition makes each step inspectable, but it guarantees neither that the steps are correct nor that the answer follows from them: chain-of-thought explanations can systematically misrepresent the true reasons for a model's prediction~\cite{turpin2023unfaithful}. For driving, this means that a structured reasoning format is not in itself evidence that a model's decision depends on its intermediate answers; the dependence has to be measured. Scoring the quality of the reasoning text, as recent driving benchmarks do~\cite{ishaq2025drivelmmo1}, does not measure it either.

\subsection{Language-model drivers}

GPT-Driver reformulates motion planning as language modelling~\cite{mao2023gptdriver}, DriveGPT4 produces interpretable end-to-end driving outputs with a multimodal large language model~\cite{xu2024drivegpt4}, Driving-with-LLMs fuses object-level vectors into the language model~\cite{chen2024drivingwithllms}, and LMDrive drives in closed loop from language instructions~\cite{shao2024lmdrive}. EMMA maps raw camera data to trajectories, objects and road-graph elements with a multimodal foundation model~\cite{hwang2025emma}. ORION guides action generation with a vision-language model~\cite{fu2025orion}, AutoVLA unifies reasoning and action tokens with fast and slow thinking modes~\cite{zhou2025autovla}, and Alpamayo-R1 couples reasoning with action prediction for long-tail driving~\cite{nvidia2025alpamayo}. SimLingo aligns language and action for vision-only closed-loop driving and won the 2024 CARLA Challenge~\cite{renz2025simlingo}. These systems build on visual instruction tuning~\cite{liu2023llava} and chain-of-thought prompting~\cite{wei2022cot}. Reinforcement learning is now also used to train the reasoning and latent representations of such models: AutoVLA fine-tunes its reasoning and actions with group relative policy optimisation~\cite{zhou2025autovla}, FLARE (co-authored by the author) activates a pre-trained vision-language model for camera-only planning on NAVSIM without language supervision by combining self-supervised future DINOv2 feature prediction with group relative policy optimisation~\cite{xie2026flare}, and on the long-tail WOD-E2E benchmark such post-training gave much larger gains when its reward was the benchmark's own metric than when it was a displacement error~\cite{xu2026wode2e}, which leaves open how much of the gain reflects better driving rather than a closer fit to the metric. Evidence that language reasoning improves driving, as opposed to explainability, is mixed and mostly open-loop; the strongest closed-loop evidence comes from systems that tightly align language with action~\cite{renz2025simlingo} or use reasoning adaptively~\cite{zhou2025autovla}.

\subsection{Dual and fast--slow systems}

Because large vision-language models are slow and weak at precise spatial reasoning, several systems pair them with a conventional planner. DriveVLM-Dual combines a vision-language model with a traditional pipeline and was deployed on a production vehicle~\cite{tian2024drivevlm}. Senna lets a vision-language model make high-level decisions and an end-to-end model produce trajectories, arguing that end-to-end models lack common sense in rare scenarios~\cite{jiang2026senna}. ETA (co-authored by the author) runs a large model one step ahead of a small, reactive one~\cite{hamdan2025eta}. How the slow and fast systems should communicate, whether through text, latent features or trajectories, and how the reasoning of the slow system can be kept grounded, are unresolved.

\subsection{Reliability and hallucination}

Object hallucination was first quantified in image captioning~\cite{rohrbach2018object}, and large vision-language models also generate objects that are inconsistent with the image~\cite{li2023pope}. CODA-LM evaluates such models automatically on driving corner cases~\cite{chen2025codalm}, and LingoQA documents a large gap to human truthfulness~\cite{marcu2024lingoqa}. The DriveBench study (co-authored by the author) showed that such models can retain much of their score even when the image is removed~\cite{xie2025vlms}. Whether a structured reasoning format constrains hallucination, or merely formats it, has not been tested systematically.

\paragraph{What is known and what is not.} Vision-language models can describe and answer questions about driving scenes, can be coupled to planners in dual systems, and can output trajectories directly~\cite{hwang2025emma,tian2024drivevlm,renz2025simlingo}, but they hallucinate and fall well short of human truthfulness on driving questions~\cite{li2023pope,marcu2024lingoqa}. What is not known is whether their reasoning can be made structured and grounded, with each answer explicitly conditioned on the perception and prediction answers it depends on, and whether such structure measurably improves planning and generalisation. This is the gap behind RQ2, which Chapter~\ref{ch:drivelm} addresses. A further open question, which Chapter~\ref{ch:drivelm} does not test, is whether the final decision of such a model actually depends on its intermediate answers~\cite{turpin2023unfaithful}.

% ==============================================================================
\section{End-to-End Autonomous Driving and Its Evaluation}\label{sec:bg:e2e}

\subsection{From ALVINN to conditional imitation}

Learning to drive from data has a long history. ALVINN learned steering with a neural network from sensor input~\cite{pomerleau1988alvinn}, and NVIDIA's convolutional network mapped raw camera pixels directly to steering commands~\cite{bojarski2016end}. Conditional imitation learning conditioned the policy on high-level navigation commands to resolve ambiguity at intersections~\cite{codevilla2018end}; ChauffeurNet added synthesised perturbations because pure cloning did not suffice~\cite{bansal2019chauffeurnet}; and Learning by Cheating first trained a privileged agent on the ground-truth scene layout and then distilled it into a vision-only student~\cite{chen2019lbc}. This early line already contains the two failure modes that recur throughout the field: ambiguity of intent, addressed by conditioning, and covariate shift, addressed by perturbation or privileged distillation.

\subsection{Differentiable modular stacks}

A second generation keeps intermediate representations while training the whole system towards planning. TransFuser fuses camera and LiDAR with transformers~\cite{chitta2023transfuser}, TCP couples trajectory planning with control prediction~\cite{wu2022tcp}, ST-P3 learns spatiotemporal features shared by perception, prediction and planning~\cite{hu2022stp3}, and VAD plans on a vectorised scene~\cite{jiang2023vad}. UniAD (co-authored by the author) chains tracking, mapping, motion forecasting, occupancy prediction and planning through learnable queries~\cite{hu2023uniad}. A survey of the field names causal confusion, robustness, interpretability and closed-loop evaluation as its open challenges~\cite{chen2024e2esurvey}. Which intermediate tasks are actually necessary is contested: a systematic study of the design space led PARA-Drive to a fully parallel architecture with nearly three times faster runtime~\cite{weng2024paradrive}, and a careful analysis on CARLA showed that reported improvements are often confounded by changes to minor system components~\cite{jaeger2023hidden}.

\subsection{Multimodal planners}

Current planners increasingly treat planning as selection among, or generation of, many candidate trajectories. Hydra-MDP distils from both human and rule-based teachers into a decoder with heads tailored to different evaluation metrics~\cite{li2024hydramdp}; it is the base planner of Chapter~\ref{ch:centaur}. DrivoR compresses multi-camera features into register tokens and uses a scoring decoder that mimics an oracle~\cite{kirby2026drivor}. Generative alternatives include truncated diffusion~\cite{liao2025diffusiondrive} and latent world models used as self-supervision~\cite{li2025law}, and reinforcement learning in photorealistic Gaussian-splatting replicas of real scenes reports substantially fewer collisions than imitation-based methods on its own benchmark~\cite{gao2025rad}. The contest between learned and rule-based planning is not settled. A rule-based planner built on a centreline prior won the 2023 nuPlan challenge~\cite{dauner2023parting}, and leading learned planners distil rule-based scores~\cite{li2024hydramdp} or mimic an oracle scorer~\cite{kirby2026drivor}; a 2026 test-time search over such a learned scorer nearly matches the privileged rule-based PDM-Closed planner on NAVSIM-v2~\cite{xu2026toad}. The frontier thus approaches rule-based performance by learning rule-derived scores, which leaves open how these planners behave where the rules do not anticipate the situation.

\subsection{Evaluation protocols}

Offline prediction metrics correlate only weakly with driving quality~\cite{codevilla2018offline}, and on nuScenes a planner that uses only the ego vehicle's state is competitive in open-loop evaluation~\cite{zhai2023rethinking,li2024ego}. Closed-loop evaluation in the CARLA simulator~\cite{dosovitskiy2017carla} was extended to multi-ability testing by Bench2Drive, whose authors argue that open-loop displacement and collision metrics cannot reflect driving performance~\cite{jia2024bench2drive}. nuPlan introduced closed-loop planning evaluation on real-world data~\cite{karnchanachari2024nuplan}. NAVSIM takes a middle ground: it unrolls bird's-eye-view abstractions of real scenes for a short horizon in a non-reactive simulation and computes rule-based metrics that align better with closed-loop results than displacement errors~\cite{dauner2024navsim}; it is built on the OpenScene redistribution of nuPlan logs described in Chapter~\ref{ch:occnet}. Pseudo-simulation further synthesises observations near the planned trajectory with 3D Gaussian splatting and raises the correlation with closed-loop simulation beyond that of open-loop metrics~\cite{cao2025pseudo}. Each protocol trades realism, reactivity and scale differently, which motivates the explicit ladder of evaluation fidelity of Chapter~\ref{ch:methodology}.

Benchmarks have also begun to target the long tail explicitly. interPlan adds edge cases and challenging scenarios to closed-loop evaluation on nuPlan, and neither the rule-based nor the learned planners it tested handled them safely~\cite{hallgarten2024interplan}. WOD-E2E, the benchmark of the 2025 Waymo vision-based end-to-end driving challenge, curates 4,021 real-world segments whose scenarios occur with a frequency below 0.03\%, and scores a trajectory by its agreement with rater-annotated preferences rather than with the single logged trajectory~\cite{xu2026wode2e}. Rater-based scoring addresses the multimodality that displacement errors ignore, and its authors report that it correlates only mildly with displacement error, but the evaluation remains open-loop. The \texttt{navsafe} split of Chapter~\ref{ch:centaur} belongs to the same family. It differs in that its frames are a subset of \texttt{navtest}, selected in part by the performance of the base planner that Chapter~\ref{ch:centaur} adapts, so it is not a held-out test set.

\paragraph{What is known and what is not.} End-to-end planners are competitive on real-data benchmarks, open-loop displacement metrics are unreliable, rule-based priors remain strong, and simulation-based metrics improve alignment with closed-loop driving~\cite{codevilla2018offline,li2024ego,dauner2023parting,dauner2024navsim,cao2025pseudo}. What is not known is how these planners behave on scenarios that are rare in the training logs, which long-tail benchmarks have only begun to measure~\cite{hallgarten2024interplan,xu2026wode2e}, and how a deployed planner can improve without retraining or hand-engineered fallbacks. Benchmarks dominated by nominal driving make safety-critical behaviour hard to measure. These gaps motivate RQ3 and Chapter~\ref{ch:centaur}.

% ==============================================================================
\section{Robustness via Adaptation}\label{sec:bg:adapt}

\subsection{Training-time domain adaptation and randomisation}

Learning theory bounds the error on a target domain by the error on the source domain plus a measure of divergence between them~\cite{bendavid2010theory}. Domain-adversarial training learns features that cannot discriminate between source and target~\cite{ganin2016dann}, domain randomisation varies simulated rendering so that the real world appears as one more variation~\cite{tobin2017domain}, and rapid motor adaptation trains a module that lets a legged-robot policy adjust online~\cite{kumar2021rma}. These methods require target data, or an anticipated family of variations, at training time; they do not cover shifts that are unknown until deployment.

\subsection{Test-time training and test-time adaptation}

Test-time training updates a model on each test input through a self-supervised auxiliary task~\cite{sun2020ttt}. TENT adapts a model by minimising the entropy of its predictions~\cite{wang2021tent}, MEMO robustifies a single test point by marginalising over augmentations~\cite{zhang2022memo}, and CoTTA targets continually changing target domains~\cite{wang2022cotta}; a recent survey organises the area~\cite{liang2025ttasurvey}. The idea has also been carried to sequential and reasoning tasks: test-time policy adaptation during deployment~\cite{hansen2021pad}, test-time training for trajectory prediction~\cite{park2024t4p}, and test-time training for few-shot reasoning in language models~\cite{akyurek2025ttt}. Unsupervised objectives such as entropy can, however, reinforce confident errors and accumulate them over time~\cite{wang2022cotta}, and most evidence remains in classification and segmentation. Test-time parameter updates for safety-critical planning, the subject of Chapter~\ref{ch:centaur}, are rare.

The failure modes of test-time adaptation on non-stationary data are by now well documented, and they bear directly on sequential deployment. Most methods assume that test samples arrive independently, whereas in applications such as driving they are temporally correlated, and under such streams most existing methods fail badly~\cite{gong2022note}. Under mixed shifts, small batches or imbalanced label streams, adaptation can fail to help or even harm, and entropy minimisation can collapse to trivial solutions that assign the same class to every input~\cite{niu2023sar}. Over long, continually changing streams, all but one of the state-of-the-art methods tested eventually collapse and perform worse than the unadapted model, and simply resetting the model to its pre-trained weights at regular intervals does as well or better~\cite{press2023rdumb}. Hyperparameters are hard to select without labels, and no method copes with all common kinds of shift~\cite{zhao2023pitfalls}. Sequential control adds a feedback loop that these benchmarks lack: an adapted policy changes the states on which it later adapts, so adaptation errors can compound like imitation errors (Section~\ref{sec:bg:shift}). Episodic adaptation, which restarts from the source weights for every input as MEMO does~\cite{zhang2022memo}, avoids accumulation at the cost of retaining nothing. Chapter~\ref{ch:centaur} evaluates adaptation on logged frames in non-reactive simulation, where this feedback cannot arise, so these failure modes lie outside what its evaluation can reveal.

\subsection{Safety filters and rule-based fallbacks}

Control-theoretic safety filters keep a system inside a safe set, for example through control barrier functions~\cite{ames2019cbf}, and a recent review unifies these approaches~\cite{hsu2024safetyfilter}. For driving, responsibility-sensitive safety formalises safe distances and proper responses~\cite{shalevshwartz2017rss}, and car-following models such as the intelligent driver model~\cite{treiber2000idm} underpin rule-based planners, including the nuPlan-winning PDM~\cite{dauner2023parting}. SafetyNet deploys a learned planner behind a rule-based fallback layer in real urban traffic~\cite{vitelli2022safetynet}. Filters offer guarantees relative to a model, but they are hand-designed, often conservative, and cannot improve from new data.

\subsection{Test-time optimisation and selection}

Diffusion planners allow guidance at inference~\cite{janner2022diffuser}; because gradient guidance needs differentiable rewards, Diffusion-ES searches with evolutionary strategies under non-differentiable rewards instead~\cite{yang2024diffusiones}. TOAD searches at test time against a learned trajectory scorer without retraining~\cite{xu2026toad}, and in manipulation, value functions re-rank the candidate actions of a generalist policy at test time~\cite{nakamoto2024vgps}. Test-time optimisation changes the output while leaving the model unchanged, whereas test-time training changes the model; which is safer, cheaper and more sample-efficient under shift has not been compared directly for planning.

\subsection{Uncertainty-aware monitoring and conformal guarantees}

Conformal prediction provides distribution-free coverage guarantees~\cite{angelopoulos2023gentle} and has been used to plan safely around predicted trajectories~\cite{lindemann2023safe}. Runtime monitors combine a fast anomaly detector with slower language-model reasoning to trigger reactive plans~\cite{sinha2024realtime}, or monitor generative policies for erratic actions and stalled task progress~\cite{agia2024unpacking}. Conformal guarantees, however, assume exchangeability, which distribution shift breaks, and monitors that raise alarms rarely feed a correction back into the policy.

\paragraph{What is known and what is not.} Models can be adapted at test time without labels~\cite{sun2020ttt,wang2021tent,zhang2022memo}, although unsupervised adaptation can degrade or collapse on long, non-stationary or temporally correlated streams~\cite{gong2022note,press2023rdumb}; rule-based filters provide dependable fallbacks~\cite{vitelli2022safetynet,hsu2024safetyfilter}, and test-time search over learned scorers is powerful~\cite{xu2026toad}. What is not known is how to make a learned planner correct itself at deployment with an objective that tolerates multimodality, is cheap enough for real time, and does not collapse, instead of switching to hand-crafted fallbacks. This gap motivates RQ3 and RQ4.

% ==============================================================================
\section{Robot Manipulation Learning}\label{sec:bg:manip}

\subsection{Expressive imitation policies and vision-language-action models}

Diffusion Policy represents visuomotor policies as conditional action diffusion and thereby captures multimodal demonstrations~\cite{chi2023diffusion}, and ACT with the ALOHA hardware predicts chunks of actions with a transformer, showing that learning lets low-cost bimanual hardware perform fine tasks~\cite{zhao2023aloha}. Chunked execution improves temporal consistency but reduces reactivity, a tension that resurfaces at deployment. Generalist policies followed: RT-1 scaled transformer policies on real-robot data~\cite{brohan2023rt1}, RT-2 expressed actions as text tokens and co-fine-tuned vision-language models on robot and web data~\cite{zitkovich2023rt2}, and Octo~\cite{team2024octo} and OpenVLA~\cite{kim2024openvla} released open generalist policies. $\pi_0$ attaches a flow-matching action expert to a vision-language model~\cite{black2024pi0}, and $\pi_{0.5}$ co-trains on heterogeneous data to perform long-horizon tasks in entirely new homes~\cite{black2025pi05}; it is the base policy of Chapter~\ref{ch:kai0}. RDT-1B scales a diffusion transformer to 1.2B parameters~\cite{liu2025rdt}, GR00T N1 offers an open humanoid foundation model with a dual-system design~\cite{nvidia2025gr00t}, and AgiBot World (co-authored by the author) provides a large-scale manipulation platform together with the GO-1 policy~\cite{bu2025agibot_iros}. Even with large-scale pre-training, however, vision-language-action policies can latch onto modality shortcuts rather than task-relevant evidence; PerturBot (co-authored by the author) diagnoses such salience, language and motor shortcuts with the offline GroundFscore metric and mitigates them via task-preserving wrist-view perturbations, decision-relevant caption enrichment and caption-grounded trajectory expansion~\cite{liu2026perturbot}. Generalisation claims are also mostly made in pick-and-place or tabletop settings; long-horizon, contact-rich, deformable tasks with hours-long reliability requirements are much less studied.

\subsection{Data at scale}

Open X-Embodiment pooled datasets across robot embodiments~\cite{oneill2024openx}, DROID collected manipulation data in the wild~\cite{khazatsky2024droid}, and BridgeData V2 provides 60,096 trajectories across 24 environments~\cite{walke2023bridgedata}. Larger datasets are of mixed quality and harder to curate~\cite{nakamoto2024vgps}, and more data does not remove the gap between the conditions of data collection and those of deployment.

\subsection{Interactive imitation and learning from deployment}

HG-DAgger lets a human expert take control when needed instead of labelling actions without being in control~\cite{kelly2019hgdagger}. Sirius frames deployment as a division of work between human and robot, learning from interventions~\cite{liu2023robotjob}, and human-in-the-loop reinforcement learning achieves precise, dexterous manipulation~\cite{luo2025hilserl}. Advantage-weighted regression turns reinforcement learning into weighted supervised learning~\cite{peng2019awr}, and $\pi^{*}_{0.6}$ conditions a vision-language-action model on binarised advantages computed from a value function, ingesting demonstrations, on-policy rollouts and expert interventions~\cite{amin2026pistar06}. A perspective paper on self-directed learning (co-authored by the author) argues that such loops should ultimately run without a human~\cite{chen2025robot}. These methods depend on costly on-robot data and human interventions, advantage estimation for long-horizon, sparsely rewarded tasks is noisy, and how to attribute credit to the stages of long tasks is open.

\subsection{Weight-space model combination}

Averaging the weights of models fine-tuned from the same initialisation improves accuracy without extra inference cost~\cite{wortsman2022soups}, ensembling zero-shot and fine-tuned weights restores robustness to distribution shift that fine-tuning otherwise reduces~\cite{wortsman2022wiseft}, task vectors can be added or negated to edit behaviour~\cite{ilharco2023task}, and merging schemes must resolve interference between task-specific models~\cite{yadav2023ties}. In robotics, policy merging has enabled fleet-level learning without centralising data~\cite{wang2024fleet}. The evidence comes mainly from vision and language classifiers; whether weight arithmetic preserves the precise, temporally coupled behaviour of visuomotor policies, as the Model Arithmetic of Chapter~\ref{ch:kai0} requires, is barely tested.

\subsection{Deformable garment manipulation and execution at deployment}

Deformable objects challenge perception, modelling and control~\cite{sanchez2018survey}. Dynamic flinging reaches high cloth coverage within a few actions~\cite{ha2021flingbot}, bimanual folding has been made efficient~\cite{avigal2022speedfolding}, and $\pi^{*}_{0.6}$ reports hours of uninterrupted laundry folding~\cite{amin2026pistar06}. At execution time, inference latency causes pauses or jerky, out-of-distribution motions at the boundaries of action chunks, which real-time chunking mitigates without retraining~\cite{black2025rtc}. Such execution effects are usually treated as engineering details rather than as a source of train--deploy distribution shift in their own right.

\paragraph{What is known and what is not.} Expressive generative policies and vision-language-action models trained on large multi-robot datasets generalise across objects and scenes~\cite{chi2023diffusion,kim2024openvla,black2025pi05}; deployment experience with interventions and advantage-based training improves reliability~\cite{luo2025hilserl,amin2026pistar06}; and weight-space combination can add robustness in vision models~\cite{wortsman2022wiseft}. What is not known is how to obtain robust, long-horizon, real-robot garment manipulation under limited resources, by treating the mismatch between data collection, training and deployment as a single distribution-shift problem, and whether model merging and stage-aware advantage signals transfer to such policies. This gap motivates the second part of RQ3, which Chapter~\ref{ch:kai0} addresses.

% ==============================================================================
\section{Benchmarks and Challenges as Research Infrastructure}\label{sec:bg:bench}

\subsection{Datasets and challenges}

KITTI established public benchmarks for driving~\cite{geiger2012kitti}; nuScenes~\cite{caesar2020nuscenes}, the Waymo Open Dataset~\cite{sun2020waymo} and Argoverse~2~\cite{wilson2021argoverse2} scaled up sensors, cities and tasks, and SemanticKITTI added point-wise semantics and scene completion~\cite{behley2019semantickitti}. Planning benchmarks followed with nuPlan~\cite{karnchanachari2024nuplan}, the CARLA leaderboard~\cite{carla2025leaderboard} and NAVSIM~\cite{dauner2024navsim}. Public challenges have repeatedly accelerated progress: the DARPA Grand Challenge was won by Stanley~\cite{thrun2006stanley}, and the ImageNet challenge organised a decade of recognition research~\cite{russakovsky2015ilsvrc}. Donoho credits the ``common task framework'', consisting of shared data, a shared metric and a held-out referee, for much of the success of predictive modelling~\cite{donoho2017fifty}. Driving continues this pattern, with public challenges on nuPlan~\cite{dauner2023parting}, NAVSIM~\cite{dauner2024navsim}, CARLA~\cite{renz2025simlingo} and occupancy prediction (Section~\ref{sec:bg:repr}); the benchmarks and challenges to which the author contributed, discussed in Chapter~\ref{ch:methodology}, belong to the same infrastructure.

\subsection{Meta-science of benchmarks}

Datasets carry recognisable biases~\cite{torralba2011unbiased}, and freshly collected test sets for ImageNet cause accuracy drops of 11--14\%~\cite{recht2019imagenet}. The benchmark lottery argues that many factors other than algorithmic superiority decide which method is perceived as better~\cite{dehghani2021benchmark}, influential benchmarks are treated as stand-ins for general capability~\cite{raji2021everything}, and dataset use concentrates on few datasets from few institutions~\cite{koch2021reduced}. Driving provides concrete cases: open-loop nuScenes metrics can be satisfied by ego status alone~\cite{zhai2023rethinking,li2024ego}, and on the CARLA leaderboard participants found that stopping a route early maximised the infraction-based score, which led to a change of the scoring formula in 2025~\cite{carla2025leaderboard}. Surveys of data-centric AI document problems in how datasets are built and used~\cite{paullada2021data,zha2025datacentric} and show that undervalued data work causes compounding ``data cascades'' in high-stakes AI~\cite{sambasivan2021everyone}. For robots, real-world evaluation of generalist policies is hard to scale and reproduce: SIMPLER narrows the gap between simulated and real evaluation~\cite{li2024simpler}, and best-practice guidelines criticise the reporting of bare success rates without the number of trials, initial conditions or success criteria~\cite{kressgazit2024robot}.

\subsection{Statistics of evaluation}

Evaluations of learned agents usually rest on few training runs and on finite sets of frames, scenarios or trials, and common reporting practice accounts for neither source of variance. In deep reinforcement learning, results vary with random seeds, environments and implementation details, so that without significance testing and standardised reporting it is hard to tell whether an improvement is real~\cite{henderson2018matters}. Agarwal et al.\ show that point estimates from a handful of runs are unreliable, and recommend interval estimates obtained by stratified bootstrap together with robust aggregates such as the interquartile mean~\cite{agarwal2021precipice}. For success rates over tens of trials, as in real-robot evaluation, the usual normal-approximation (Wald) interval has erratic coverage, and the Wilson or Jeffreys interval is recommended for samples of up to about 40~\cite{brown2001interval}. When two methods are scored on the same frames, as in open-loop nuScenes and NAVSIM evaluation, their per-frame outcomes are paired, and a paired comparison, such as a bootstrap over frames of the per-frame difference, is usually more sensitive than a comparison of two separately averaged scores. These tools delimit what single-run results, which are common in this thesis, can and cannot show; the statistical practice adopted in the thesis is described in Section~\ref{sec:method:evaluation}.

\paragraph{What is known and what is not.} Shared datasets and challenges accelerate progress~\cite{russakovsky2015ilsvrc,donoho2017fifty}, and benchmarks can equally be overfitted, gamed and mistaken for general capability~\cite{recht2019imagenet,dehghani2021benchmark,raji2021everything}; interval estimates and paired comparisons are the minimum needed to interpret small differences~\cite{agarwal2021precipice,brown2001interval}. What is not known is how to design benchmarks for physical agents that measure behaviour under shift and in the long tail rather than on nominal logs, that co-evolve with methods before they saturate, and that standardise real-robot evaluation. This question underlies the methodology of Chapter~\ref{ch:methodology}.

% ==============================================================================
\section{Summary: Gaps and Research Questions}\label{sec:bg:summary}

Table~\ref{tab:bg:gaps} collects the gaps identified in this chapter and relates each to the research questions of Chapter~\ref{ch:intro} and to the chapter that addresses it. Two observations connect them. First, the gaps are distributed along the whole perceive--reason--act loop, which is why the thesis addresses each stage of the loop rather than one. Second, many gaps concern measurement as much as method: open-loop metrics that mislead, question-answering benchmarks without explicit structure, planning benchmarks dominated by nominal driving, and real-robot evaluation without standard protocols. This is why each core chapter contributes a benchmark or an evaluation protocol along with its method, following the methodology described in the next chapter.

\begin{table}[t]
  \centering
  \caption[Gaps in the literature and where this thesis addresses them]{Gaps in the literature identified in this chapter, the research questions they motivate, and the chapters that address them.}
  \label{tab:bg:gaps}
  \footnotesize
  \renewcommand{\arraystretch}{1.2}
  \begin{tabularx}{\textwidth}{@{}X p{1.2cm} p{1.6cm}@{}}
    \toprule
    Gap & RQ & Chapter \\
    \midrule
    Whether occupancy makes downstream planning better than boxes or vectors, and whether one representation can serve many tasks at scale (Section~\ref{sec:bg:repr}) & RQ1 & \ref{ch:occnet} \\
    Driving question answering without explicit reasoning structure; unclear benefit of language reasoning for driving under shift (Section~\ref{sec:bg:lang}) & RQ2 & \ref{ch:drivelm} \\
    Uncertainty measured but not used as a learning signal for multimodal planners (Sections~\ref{sec:bg:shift} and~\ref{sec:bg:adapt}) & RQ3, RQ4 & \ref{ch:centaur} \\
    Learned planners rely on hand-designed fallbacks rather than self-correction; safety-critical behaviour hard to measure (Sections~\ref{sec:bg:e2e} and~\ref{sec:bg:adapt}) & RQ3, RQ4 & \ref{ch:centaur} \\
    Corrective signals at deployment require humans; train--deploy mismatch in execution treated piecemeal (Section~\ref{sec:bg:manip}) & RQ3, RQ4 & \ref{ch:kai0} \\
    Model merging and stage-aware advantages untested for visuomotor policies (Section~\ref{sec:bg:manip}) & RQ3 & \ref{ch:kai0} \\
    Benchmarks that measure the long tail, co-evolve with methods and standardise real-robot evaluation (Section~\ref{sec:bg:bench}) & all & \ref{ch:methodology}, \ref{ch:discussion} \\
    \bottomrule
  \end{tabularx}
\end{table}
